% !TeX root = ../Thesis.tex

%************************************************
\chapter{Grundlagen}\label{ch:grundlagen}
%************************************************
\glsresetall % Resets all acronyms to not used

%**********Grundlagen der Numismatik*************
\section{Grundlagen der Numismatik}
\label{sec:glnumismatik}
Die Numismatik, auch Münzkunde genannt, beschäftigt sich mit der wissenschaftlichen Analyse von Münzen und dem Geldwesen. Sie berührt dabei unter anderem die Archäologie sowie die Geschichts-, Wirtschafts- und Sozialwissenschaft \cite{wikipedia-numismatik}. Im Mittelpunkt dieser Arbeit stehen antike Münzen, genauer ihre Münzbilder, als Untersuchungsgegenstand.

Aus vielen Epochen und Regionen sind Münzfunde und Sammlungen erhalten, so stammen die ältesten Münzen aus der Zeit um 600 v. Chr. \cite{vonkaenel2012ersten}. Als historische Quellen geben Münzen Auskunft über die Zeit, in der sie geprägt und im Handel genutzt wurden. Daraus lassen sich somit politische und wirtschaftliche Entwicklungen nachvollziehen. Da Münzen in großer Zahl erhalten sind, sind sie besonders für Epochen von Interesse, aus denen wenig schriftliche Überlieferung oder wenige genaue Datierungen vorliegen. Ihre Prägungen zeigen Datierungshinweise und Inschriften sowie Darstellungen von Personen und Symbolen, die auf Zeit und Kultur der Prägung verweisen. Auf antiken Münzen sind etwa Gottheiten und Heroen wie Herakles, Aphrodite oder Athena abgebildet. Die Römer übernahmen viele Motive von den Griechen, die bereits Personen darstellten. Eine Besonderheit römischer Münzen sind Abbildungen von Herrschern; Caesar war der erste Römer, der zu Lebzeiten auf einer Münze abgebildet wurde \cite{howgego1995ancient,museumdigital-roemische-muenzen} \todo{Formulierung und Seite prüfen}. 

Antike Münzen werden in Vorder- und Rückseite unterteilt, die als Obverse (auch Avers) und Reverse (auch Revers) bezeichnet werden. Bei römischen Kaisermünzen zeigt die Vorderseite in der Regel die Büste des Herrschers, die Rückseite dagegen Verschiedenes, beispielsweise das Bild einer Gottheit, Insignien, Tiere oder Gegenstände \cite{museumdigital-roemische-muenzen}. Zusätzlich trägt die Vorderseite häufig eine Legende, also eine Inschrift mit Namen, Titeln oder weiteren Angaben \cite{cooper2020learning}, siehe \cref{fig:coin-3545}. Andere antike Prägungen weichen davon ab. Frühe thrakische Münzen zeigen auf der Vorderseite häufig ein Tier oder einen Gegenstand und auf der Rückseite ein vertieftes Viereck (Quadratum Incusum) \cite{berthold2016thrakien}. Welche Motive im Datensatz dieser Arbeit vorkommen, wird in \cref{sec:gesamtansatzdatengrundlage} beschrieben. Diese sichtbaren und inhaltlich interpretierbaren Bestandteile werden in dieser Arbeit als visuelle Merkmale bezeichnet (vgl. \cref{sec:merkmale}).

\begin{figure}[H]
    \centering

    \begin{subfigure}{0.4\textwidth}
        \centering
        \includegraphics[width=\textwidth]{coin3545_obverse.jpeg}
        \caption{Vorderseite}
    \end{subfigure}
    \hspace{0.04\textwidth}
    \begin{subfigure}{0.4\textwidth}
        \centering
        \includegraphics[width=\textwidth]{coin3545_reverse.jpeg}
        \caption{Rückseite}
    \end{subfigure}

    \caption[Beispielmünze aus dem Corpus Nummorum] {Beispielmünze aus dem Corpus Nummorum. Ca. 222-235 n. Chr. Münzstätte Byzantion, Region Thrakien.
  Quelle: cn coin 3545 \cite{cn-coin-3545}. Original, Münzkabinett, Staatliche Museen
  zu Berlin, Inv.-Nr.~18235544. Foto: Reinhard Saczewski. Lizenz: Public Domain Mark 1.0.}
    \label{fig:coin-3545}
\end{figure}

\noindent Auch die Herstellung beeinflusst das Erscheinungsbild. Münzen wurden von der Antike bis in die frühe Neuzeit durch Hammerprägung hergestellt \cite{wikipedia-muenzpraegung}. Ein Münzrohling, der sogenannte Schrötling, wurde auf einen fest installierten Unterstempel gelegt, der das Motiv der Vorderseite trug. Ein frei geführter Oberstempel mit dem Motiv der Rückseite wurde daraufgesetzt und mit einem Hammerschlag eingeprägt, sodass beide Seiten geprägt wurden \cite{adfontes-herstellung}. Da jeder Stempel ein in Handarbeit gefertigtes Einzelstück war, unterscheiden sich auch Münzen desselben Typs in feinen Details \cite{wikipedia-muenzpraegung}. Außerdem nutzten sich die Stempel ab und mussten ersetzt werden, wobei der Oberstempel in der Regel früher ausgetauscht wurde als der Unterstempel. Dadurch ergeben sich unterschiedliche Stempelkombinationen, anhand derer sich die Reihenfolge der Prägung rekonstruieren lässt (Stempelanalyse) \cite{wikipedia-numismatik}.

Zusätzlich verändert sich das Erscheinungsbild im Laufe der Objektgeschichte \cite{wikipedia-numismatik}. Nutzung und Umlauf führen zu Abnutzung, die je nach Erhaltungsgrad unterschiedlich stark ausfällt \cite{wikipedia-erhaltungsgrade,adfontes-abnutzung} (\cref{fig:erhaltung}). Bei Münzen, die lange im Boden lagen, kann je nach Bodenbedingungen außerdem das Metall korrodieren. Motive, Symbole und Legenden können dadurch teilweise oder vollständig unkenntlich werden \cite{adfontes-abnutzung}. Abnutzung und Korrosion erschweren zudem die automatisierte Erkennung von Merkmalen auf antiken Münzbildern, da Details und Legenden verloren gehen \cite{fare2017ancient,heinecke2021unsupervised} (vgl. \cref{sec:modelle}).

\begin{figure}[H]
    \centering

    \begin{subfigure}{0.4\textwidth}
        \centering
        \includegraphics[width=\textwidth]{coin1018_obverse.jpeg}
        \caption{Original}
        \label{fig:originalcoin1018}
    \end{subfigure}
    \hspace{0.04\textwidth}
    \begin{subfigure}{0.4\textwidth}
        \centering
        \includegraphics[width=\textwidth]{coin1019_obverse.jpeg}
        \caption{Original}
        \label{fig:originalcoin1019}
    \end{subfigure}

    \vspace{0.6em}

    \begin{subfigure}{0.4\textwidth}
        \centering
        \includegraphics[width=\textwidth]{coin1015_obverse.jpeg}
        \caption{Gipsabguss}
        \label{fig:gibsabgusscoin1015}
    \end{subfigure}
    \hspace{0.04\textwidth}
    \begin{subfigure}{0.4\textwidth}
        \centering
        \includegraphics[width=\textwidth]{coin1017_obverse.jpeg}
        \caption{Gipsabguss}
        \label{fig:gipsabgusscoin1017}
    \end{subfigure}
    \caption[Münzen in unterschiedlichem Erhaltungszustand]{Münzen aus dem Corpus Nummorum in unterschiedlichem Erhaltungszustand, alle aus Byzantion, Region Thrakien, ca. 222--235 n.~Chr.: Originale (oben) und Gipsabgüsse (unten).
    (a) cn coin 1019~\cite{cn-coin-1019}. Original, Münzkabinett, Staatliche Museen zu Berlin, Inv.-Nr.~18235900. Foto: Reinhard Saczewski. Lizenz: Public Domain Mark 1.0.
    (b) cn coin 1018~\cite{cn-coin-1018}. Original, Bibliothèque nationale de France, Département des Monnaies, médailles et antiques, Paris, Inv.-Nr.~466. Lizenz: CC BY-NC-SA 4.0.
    (c) cn coin 1017~\cite{cn-coin-1017}. Gipsabguss. Foto: BBAW. Lizenz: CC BY-NC-SA 4.0.
    (d) cn coin 1015~\cite{cn-coin-1015}. Gipsabguss. Foto: BBAW. Lizenz: CC BY-NC-SA 4.0.}
    \label{fig:erhaltung}
\end{figure}

%**********Grundlagen des maschinellen Lernens*************
\section{Grundlagen des maschinellen Lernens}
\label{sec:glmaschinelleslernen}
Die technische Grundlage dieser Arbeit bildet das maschinelle Lernen, ein Teilgebiet der künstlichen Intelligenz. Im Gegensatz zu regelbasierten Verfahren wird der Lösungsweg nicht vollständig vorgegeben. Ein Modell wird anhand von Beispieldaten trainiert und kann anschließend auch für unbekannte Daten Vorhersagen treffen \cite{goodfellow2016deep,lecun2015deep} \todo{Goodfellow: Kap.5 (Learning Algorithms, Generalisierung) Seiten?}. 
Beim überwachten Lernen (Supervised Learning) stehen im Training zu den Eingabedaten die erwarteten Ausgaben zur Verfügung, etwa Bilder mit zugeordneter Klasse. Beim unüberwachten Lernen fehlen solche Zielwerte \cite{goodfellow2016deep,lecun2015deep}. Da Objektdetektoren in der Regel überwacht trainiert werden (vgl. \cref{sec:ooyolo}), ist für diese Arbeit das überwachte Lernen maßgeblich. 

Eine Weiterentwicklung stellt das Deep Learning dar, dessen Grundlage künstliche neuronale Netze sind. Ein solches Netz besteht aus mehreren hintereinandergeschalteten Schichten, nämlich: einer Eingabeschicht, einer oder mehreren verdeckten Schichten (Hidden Layers) und einer Ausgabeschicht \cref{fig:nn}. 

\begin{figure}[H]
  \centering
  \resizebox{0.7\linewidth}{!}{\begin{tikzpicture}[x=1cm,y=1cm]
  % Netz 4-4-4-2 (nur Aufbau: Schichten und Verbindungen)
  \foreach \l/\n/\col in {0/4/figblue!20,1/4/figorange!20,2/4/figorange!20,3/2/figgreen!25} {
    \foreach \i in {1,...,\n} {
      \node[figneuron, fill=\col] (n\l-\i) at (\l*1.5, {(\n+1)/2*0.8-\i*0.8}) {};
    }
  }
  \foreach \l/\m/\cnt in {0/1/4,1/2/4,2/3/2} {
    \foreach \i in {1,...,4} {
      \foreach \j in {1,...,\cnt} {
        \draw[black!35, line width=0.3pt] (n\l-\i) -- (n\m-\j);
      }
    }
  }
  \node[figlabel] at (0,-2.4) {Eingabe};
  \node[figlabel] at (2.25,-2.4) {verdeckte Schichten};
  \node[figlabel] at (4.5,-2.4) {Ausgabe};
\end{tikzpicture}
}
  \caption[Aufbau eines neuronalen Netzes]{Aufbau eines künstlichen neuronalen Netzes aus Eingabeschicht, verdeckten Schichten und Ausgabeschicht. Eigene Darstellung nach dem Prinzip aus \cite{goodfellow2016deep}.}
  \label{fig:nn}
\end{figure}

Verdeckt heißen die mittleren Schichten, weil die Trainingsdaten nur die gewünschte Ausgabe vorgeben, nicht aber, was in diesen Schichten berechnet werden soll \cite{goodfellow2016deep}. Jede Schicht erhält die Werte der vorherigen Schicht, verrechnet sie mit gewichteten Verbindungen und reicht das Ergebnis weiter. Die Gewichte bestimmen, wie stark ein Eingangswert in das Ergebnis eingeht. Durch eine Aktivierungsfunktion, in modernen Netzen meist die ReLU-Funktion, kann das Netz auch nichtlineare Zusammenhänge abbilden. Ohne sie ließe sich selbst ein Netz mit mehreren Schichten auf eine einzige lineare Abbildung zurückführen \cite{goodfellow2016deep, lecun2015deep}. Die Gewichte und die weiteren einstellbaren Größen des Netzes heißen zusammen Parameter. Sie werden im Training schrittweise angepasst. Die Ausgabe des Netzes wird mit dem erwarteten Ergebnis verglichen, und eine Kostenfunktion (Verlustfunktion, Loss Function) misst die Abweichung. Mit der Backpropagation wird berechnet, wie stark jeder Parameter zu diesem Fehler beiträgt. Dazu wird der Fehler mithilfe der Kettenregel (Kettenregel der Differentialrechnung) rückwärts durch das Netz geführt. Anhand dieser Werte werden die Parameter dann so verändert, dass der Fehler sinkt. Über viele Durchläufe durch die Trainingsdaten lernt das Netz so, welche Eigenschaften der Eingabe für die Aufgabe wichtig sind \cite{goodfellow2016deep, lecun2015deep}. Der Ablauf ist in \cref{fig:nn-training} dargestellt.

\begin{figure}[H]
  \centering
  \resizebox{\linewidth}{!}{% Optional: nur einbinden, wenn der Text Verlustfunktion und Backpropagation beschreibt
\begin{tikzpicture}
  \node[figbox, sblue] (in) at (0,0) {Eingabe};
  \node[figbox, sorange] (net) at (3,0) {Netz};
  \node[figbox, sblue] (out) at (6,0) {Ausgabe};
  \node[figbox, sorange] (loss) at (9,0) {Verlust-\\funktion};
  \node[figbox, sgreen] (exp) at (9,-1.8) {erwartetes\\Ergebnis};
  \draw[figarrow] (in) -- (net);
  \draw[figarrow] (net) -- (out);
  \draw[figarrow] (out) -- (loss);
  \draw[figarrow] (exp) -- (loss);
  \draw[figarrow, dashed] (loss.north) -- ++(0,0.8) -| (net.north)
    node[pos=0.25, above, figlabel] {Backpropagation};
\end{tikzpicture}
}
  \caption[Training eines neuronalen Netzes]{Trainingsablauf eines neuronalen Netzes: Die Ausgabe des Netzes wird mit dem erwarteten Ergebnis verglichen, die Verlustfunktion misst die Abweichung, und per Backpropagation werden die Parameter des Netzes angepasst. Eigene Darstellung nach dem Prinzip aus \cite{goodfellow2016deep,lecun2015deep}.}
  \label{fig:nn-training}
\end{figure}

\noindent Für Bilddaten haben sich Convolutional Neural Networks (CNNs) etabliert, deren Architektur auf räumlich strukturierte Daten ausgelegt ist \cite{goodfellow2016deep, lecun2015deep}. Ihr zentrales Element sind Faltungsschichten (Convolutional Layers). Sie verwenden Filter, kleine Gewichtsmatrizen, die über das Bild geschoben werden und dabei lokale Bildbereiche auf bestimmte Strukturen untersuchen (\cref{fig:faltung}). Derselbe Filter wird dabei an allen Bildpositionen verwendet, sodass das Netz weniger Parameter benötigt \cite{goodfellow2016deep}. 

\begin{figure}[H]
  \centering
  \resizebox{\linewidth}{!}{\begin{tikzpicture}[x=0.6cm,y=0.6cm, every node/.style={font=\small}]
  % Eingabe 5x5: Zeilen "0 0 9 9 9"
  \foreach \r in {0,...,4}{
    \foreach \c/\v/\g in {0/0/0,1/0/0,2/9/27,3/9/27,4/9/27}{
      \node[draw=black!50, minimum size=0.6cm, inner sep=0pt, fill=black!\g] at (\c,-\r) {\v};
    }
  }
  % Hervorgehobener Ausschnitt (Zeilen/Spalten 1-3 -> Index 0..2)
  \draw[figblue, line width=1.6pt] (-0.5,0.5) rectangle (2.5,-2.5);
  \node[figlabel] at (2,1.2) {Bildausschnitt (5$\times$5)};

  \node at (5.1,-1) {$\ast$};
  % Kernel
  \foreach \r in {0,1,2}{
    \foreach \c/\v in {0/1,1/0,2/-1}{
      \node[draw=figorange, fill=figorange!15, minimum size=0.6cm, inner sep=0pt] at (6.2+\c,-0.0-\r-0) {\v};
    }
  }
  \node[figlabel] at (7.2,1.2) {Filter (3$\times$3)};
  \node at (10.3,-1) {$=$};
  % Ausgabe 3x3: je Zeile -27 -27 0
  \foreach \r in {0,1,2}{
    \foreach \c/\v in {0/-27,1/-27,2/0}{
      \node[draw=black!50, minimum size=0.6cm, inner sep=0pt, fill=figgreen!15] at (11.4+\c,-\r) {\v};
    }
  }
  \draw[figblue, line width=1.6pt] (10.9,0.5) rectangle (11.9,-0.5);
  \node[figlabel] at (12.4,1.2) {Ergebnis (3$\times$3)};
  \node[figlabel, text width=7cm] at (6.5,-5.6) {Beispiel: $3\cdot(0\cdot1+0\cdot0+9\cdot(-1)) = -27$\\(großer Betrag bei einer senkrechten Kante)};
\end{tikzpicture}
}
  \caption[Faltung mit einem Filter]{Funktionsprinzip eines Filters in einer Faltungsschicht: Der Filter wird über den Bildausschnitt geschoben, und an jeder Position werden die Produkte aus Bild- und Filterwerten summiert. Der markierte Ausschnitt ergibt $3\cdot(0\cdot1+0\cdot0+9\cdot(-1))=-27$, der Betrag ist bei der senkrechten Kante groß. Eigene Darstellung mit beispielhaft gewählten Zahlenwerten.}
  \label{fig:faltung}
\end{figure}

\noindent Frühe Schichten erfassen einfache Strukturen wie Kanten, spätere setzen diese zu komplexeren Mustern und Formen zusammen. So entsteht eine hierarchische Repräsentation des Bildes. Anders als bei klassischen Verfahren der Bildverarbeitung müssen die relevanten Merkmale dabei nicht manuell festgelegt werden, sondern werden aus den Trainingsdaten erlernt \cite{lecun2015deep, goodfellow2016deep}. Pooling-Schichten fassen benachbarte Werte zusammen, beispielsweise über den Maximalwert oder den Durchschnitt, und verkleinern so die räumliche Auflösung der Merkmalskarten \cite{goodfellow2016deep}. 

\noindent Ein Modell muss nicht für jede Aufgabe von Grund auf trainiert werden. Beim Transfer Learning werden Repräsentationen, die ein Modell auf einem großen, allgemeinen Datensatz gelernt hat, für eine andere Zielaufgabe weiterverwendet. Wird das vortrainierte Modell anschließend mit gelabelten Daten der Zielaufgabe weiter trainiert, spricht man von Fine-Tuning. Beim Few-Shot Learning genügen dafür wenige Beispiele der Zielaufgabe, beim Zero-Shot Learning kommt das Modell ganz ohne solche Beispiele aus. Stattdessen wird zusätzliches Wissen über die Zielklassen benötigt, etwa in Form von Attributen oder Textbeschreibungen \cite{zhang2019recent}. Ein Beispiel sind Modelle, die
Bilder anhand von Texten ihrer Klassen einordnen \cite{zhai2022lit}.
Dieses Prinzip nutzt das in dieser Arbeit eingesetzte YOLO-World \cite{cheng2024yoloworld} (vgl. \cref{sec:modelle}). Es gehört zur Objekterkennung, einer zentralen Aufgabe der Computer Vision, die im folgenden Abschnitt behandelt wird. 

%**********Objekterkennung und YOLO*************
\section{Objekterkennung und YOLO}
\label{sec:ooyolo}
Die Objekterkennung (Object Detection), im Folgenden synonym als Detektion bezeichnet, ist eine zentrale Teilaufgabe der Computer Vision \cite{redmon2016yolo}. Im Gegensatz zur Bildklassifikation, die einem gesamten Bild eine einzelne Klasse zuordnet, wird bei der Detektion für jedes Objekt zusätzlich seine Position bestimmt (\cref{fig:klassifikation-detektion}). Das Ergebnis ist je Objekt eine Bounding Box, ein rechteckiger Bildausschnitt, der Position und Ausdehnung markiert, sowie eine Klassenzuordnung. Dadurch lassen sich mehrere Objekte in einem Bild unabhängig voneinander lokalisieren \cite{ultralytics-detection}.

\begin{figure}[H]
  \centering
  \resizebox{\linewidth}{!}{\begin{tikzpicture}
  \node[inner sep=0] (a) at (0,0) {\coinpic{3.4cm}};
  \node[figbox, sblue, below=3mm of a] {eine Klasse für das ganze Bild};
  \node[figlabel, above=2mm of a] {Klassifikation};
  \node[inner sep=0] (b) at (6,0) {\coinpic{3.4cm}};
  \node[figlabel, above=2mm of b] {Detektion};
  \begin{scope}[imgcoords=b, line width=1.2pt]
    \draw[figgreen] (0.30,0.45) rectangle (0.72,0.92) node[above left, font=\scriptsize, fill=white, inner sep=1pt] {Porträt};
    \draw[figorange] (0.08,0.05) rectangle (0.92,0.22) node[below right, font=\scriptsize, fill=white, inner sep=1pt] {Buchstabe};
    \draw[figpurple] (0.72,0.30) rectangle (0.95,0.50) node[above right, font=\scriptsize, fill=white, inner sep=1pt] {Symbol};
  \end{scope}
  \node[figbox, sgreen, below=3mm of b] {Bounding Box + Klasse je Objekt};
\end{tikzpicture}
}
  \caption[Klassifikation und Objekterkennung]{Bildklassifikation ordnet dem gesamten Bild eine Klasse zu, die Objekterkennung liefert für jedes Objekt eine Bounding Box und eine Klasse. Schematische eigene Darstellung; Münzbild: \cref{fig:coin-3545}.}
  \label{fig:klassifikation-detektion}
\end{figure}

\noindent Für das überwachte Training eines Detektors (vgl. \cref{sec:glmaschinelleslernen}) müssen die Trainingsbilder annotiert sein. Zu jedem Objekt werden Klasse und Position, in der Regel in Form einer Bounding Box angegeben. Ein so trainierter Detektor kann ausschließlich die Klassen detektieren, die in den Trainingsdaten annotiert waren \cite{cheng2024yoloworld}. Zweistufige Detektoren wie R-CNN \cite{girshick2014rich} schlagen zunächst Bildregionen vor, die ein Objekt enthalten könnten, und klassifizieren diese anschließend einzeln \cite{redmon2016yolo,cheng2024yoloworld}. Da pro Bild mehrere Verarbeitungsschritte nötig sind, ist dieses Vorgehen rechenintensiv \cite{redmon2016yolo}. Das YOLO-Verfahren (You Only Look Once) wurde 2016 von Redmon et al. \cite{redmon2016yolo} vorgestellt und verfolgt einen anderen Ansatz. Es formuliert die Detektion als einzelnes Regressionsproblem. Ein neuronales Netz sagt Koordinaten der Bounding Boxes und Klassenwahrscheinlichkeiten als Zahlenwerte direkt aus dem Bild vorher, und zwar in einem einzigen Durchlauf. Daraus leitet sich der Name ab, denn das Bild wird dem Netz nur einmal gezeigt. Bereits das ursprüngliche Modell arbeitet in Echtzeit; Redmon et al. geben 45 Bilder pro Sekunde an \cite{redmon2016yolo}.

Zur Veranschaulichung des Grundprinzips wird zunächst die ursprüngliche Architektur (YOLOv1) beschrieben. Dort wird das Eingabebild in ein regelmäßiges Raster aus Zellen unterteilt. Jede Zelle ist für Objekte zuständig, deren Mittelpunkt in ihr liegt. Sie sagt eine feste Anzahl von Bounding Boxes vor, jeweils mit Koordinaten und einem Confidence Score. Dieser Wert gibt an, wie sicher sich in der Box ein Objekt befindet und wie gut die Box dessen tatsächliche Position trifft. Zusätzlich bestimmt die Zelle, zu welcher Klasse ein enthaltenes Objekt am wahrscheinlichsten gehört \cite{redmon2016yolo} ()\cref{fig:yolov1raster}. 

\begin{figure}[H]
    \centering
    \begin{subfigure}{0.9\textwidth}
        \centering
        \includegraphics[width=\textwidth]{yolov1raster.jpg}
    \end{subfigure}
    \caption[YOLOv1 Raster] {Vereinfachte YOLOv1 Ausgabevorhersage bei einem 9x9 Raster mit 3 Klassen (blau, grün, gelb).
  Quelle: \cite{terven2023review}, Abb.~5, CC BY 3.0, ins Deutsche übertragen.}
    \label{fig:yolov1raster}
\end{figure}

\noindent Spätere Generationen wie YOLOv8 \cite{jocher2023yolov8} haben dieses Prinzip weiterentwickelt, unter anderem durch Vorhersagen auf mehreren Skalen und einen entkoppelten Head \cite{cheng2024yoloworld}. Es ist daher nicht unverändert auf moderne YOLO-Modelle übertragbar. Moderne YOLO-Netze lassen sich vereinfacht in drei Abschnitte gliedern \cite{cheng2024yoloworld} ()\cref{fig:backbone-neck-head}. Der Backbone extrahiert aus dem Eingabebild Features und überführt es schrittweise in eine kompakte Repräsentation. Der Neck führt diese Repräsentationen aus verschiedenen Ebenen zusammen und bereitet sie auf. Der Head berechnet daraus die finalen Bounding Boxes, Confidence Scores und Klassenwahrscheinlichkeiten. 

\begin{figure}[H]
  \centering
  \resizebox{\linewidth}{!}{\begin{tikzpicture}[node distance=7mm]
  \node[figbox] (in) {Eingabe-\\bild};
  \node[figbox, sblue, right=of in, minimum height=2.2cm, minimum width=1.1cm] (b1) {};
  \node[figbox, sblue, right=1mm of b1, minimum height=1.7cm, minimum width=0.9cm] (b2) {};
  \node[figbox, sblue, right=1mm of b2, minimum height=1.3cm, minimum width=0.7cm] (b3) {};
  \node[figbox, sblue, right=1mm of b3, minimum height=0.9cm, minimum width=0.5cm] (b4) {};
  \node[figbox, sorange, right=14mm of b4, minimum height=1.8cm, text width=2.2cm] (neck) {Neck\\[-1pt]{\footnotesize führt Merkmale zusammen}};
  \node[figbox, sgreen, right=of neck, minimum height=1.8cm, text width=2.6cm] (head) {Head\\[-1pt]{\footnotesize Bounding Boxes, Confidence Scores, Klassen}};
  \node[figbox, right=of head] (out) {Aus-\\gabe};
  \node[figlabel, above=2mm of b2.north east] {Backbone};
  \draw[figarrow] (in) -- (b1);
  \draw[figarrow] (b4) -- (neck);
  \draw[figarrow] (neck) -- (head);
  \draw[figarrow] (head) -- (out);
\end{tikzpicture}
}
  \caption[Aufbau eines YOLO-Netzwerks] {Vereinfachter Aufbau eines YOLO-Netzwerks aus Backbone, Neck und Head. Eigene Darstellung nach dem Aufbau aus \cite{cheng2024yoloworld}.}
  \label{fig:backbone-neck-head}
\end{figure}

\noindent Die in Backbone und Neck entstehenden Repräsentationen entsprechen den Features im Sinne von \cref{sec:merkmale}. Die in dieser Arbeit eingesetzten Modelle bauen auf YOLOv8 \cite{cheng2024yoloworld} auf (vgl. \cref{sec:modelle}).

%**********Merkmale im Kontext dieser Arbeit*************
\section{Merkmale im Kontext dieser Arbeit}
\label{sec:merkmale}
Der Begriff „Merkmal“, im Englischen „Feature“, wird je nach Fachgebiet unterschiedlich verwendet. Für die Detektion und den Vergleich von Münzbildern sind zwei Bedeutungen relevant, die im Folgenden voneinander abgegrenzt werden.
Im numismatischen Sprachgebrauch bezeichnet ein Merkmal ein für Menschen sichtbares und inhaltlich interpretierbares Element einer Münze, beispielsweise ein Herrscherporträt, eine Legende, ein Symbol oder eine Tierdarstellung (vgl. \cref{sec:glnumismatik}). Ein solches Merkmal lässt sich räumlich eingrenzen und einer Kategorie zuordnen. In diesem Sinne wird der Begriff in dieser Arbeit verwendet: Ein visuelles Merkmal ist ein abgrenzbarer Bildbestandteil, der von YOLO als Objekt lokalisiert und einer Klasse zugeordnet wird (vgl. \cref{sec:ooyolo}).
In der Bildverarbeitung und im maschinellen Lernen bezeichnet ein Feature dagegen eine numerische Eigenschaft der Eingabedaten, anhand derer ein Modell Bilddaten beschreibt und unterscheidet. Neuronale Netze erlernen solche Features während des Trainings selbstständig (vgl. \cref{sec:glmaschinelleslernen}). Mehrere Features werden zu einem Feature-Vektor 
$ A = (a_1, a_2, \dots, a_n) $ zusammengefasst, einer geordneten Folge von Werten \cite{duda2001pattern,bishop2006pattern} \todo{Bücher: Kapitel/Seiten nachschauen}. Die einzelnen Komponenten sind für Menschen in der Regel nicht interpretierbar \cite{montavon2018methods}. 
Für den Bildausschnitt eines detektierten Merkmals lässt sich ein Feature-Vektor aus einer internen Schicht des Netzwerks bestimmen \cite{ultralytics-predict,zhang2019recent}. Das Merkmal ist damit die für Menschen interpretierbare Ebene, der Feature-Vektor die numerische Repräsentation desselben Bildausschnitts. Die Extraktion wird in \cref{sec:vorgehen} beschrieben, die technische Umsetzung in \cref{sec:implementierung}.
Zwei Feature-Vektoren $ A $ und $ B $ gleicher Dimension $ n $ lassen sich mithilfe von Distanz- oder Ähnlichkeitsmaßen vergleichen. In dieser Arbeit werden drei gebräuchliche Maße verwendet:
\begin{itemize}
    \item Die Manhattan-Distanz summiert die Beträge der Differenzen \cite{aggarwal2001surprising}: \[ d_1(A, B) = \sum_{i=1}^{n}|a_i - b_i| \] Abweichungen gehen linear ein, einzelne große Komponenten dominieren das Ergebnis daher weniger.
    \item Die euklidische Distanz ist der geradlinige Abstand zweier Punkte im $n$-dimensionalen Raum \cite{wikipedia-euklidischer-abstand}: \[ d_2(A, B)= \sqrt{\sum_{i=1}^{n}(a_i - b_i)^2} \] Da die Differenzen quadriert werden, gehen einzelne große Abweichungen besonders stark in den Wert ein.
    \item Die Kosinus-Ähnlichkeit misst den Winkel zwischen zwei Vektoren \cite{manning2008introduction}: \[ \cos(A, B) = \frac{A \cdot B}{||A|| \cdot ||B||} \] Der Wert liegt zwischen $ -1 $ und $ 1 $. Er hängt nur von der Richtung der Vektoren ab, nicht von ihrer Länge.
    \todo{Buch: Abschnitt 6.3 bzw. Goodfellow nachschauen}
\end{itemize}
Bei den beiden Distanzen bedeutet ein kleiner Wert eine hohe Ähnlichkeit, bei der Kosinus-Ähnlichkeit ein großer. Die drei Maße gewichten Unterschiede zwischen Vektoren verschieden. Welches sich für den Vergleich von Merkmalen auf Münzbildern eignet, ist vorab nicht bekannt. Die Auswahl und Anwendung werden ebenfalls in \cref{sec:vorgehen} beschrieben.